% Epilogo: fora das Partes, sem numero, mas no sumario.
\chapter*{Epílogo}
\addcontentsline{toc}{part}{Epílogo}
\markboth{Epílogo}{Epílogo}
\label{cap:epilogo}

Cada onda contada neste livro prometeu, de alguma forma, acabar com o trabalho de quem analisa
dados. A planilha ia dispensar o centro de processamento. O data warehouse ia acabar com as
extrações manuais. O self-service ia dispensar os analistas, porque cada gestor montaria o
próprio painel. O big data ia encontrar as respostas sozinho, sem que ninguém precisasse fazer
perguntas. E agora a IA vai escrever as consultas, montar os modelos e, no estágio proativo,
tomar as decisões.

Nenhuma dessas promessas se cumpriu do jeito que foi feita, e todas mudaram o trabalho de
verdade. A planilha levou a análise para perto de quem decide. O data warehouse trouxe a ideia de
uma base única e integrada. O self-service encurtou a distância entre a pergunta e o gráfico. A
nuvem tornou acessível o que antes exigia um projeto de anos. Cada onda tirou do caminho uma parte
do trabalho e deixou mais visível a parte que restava.

A IA está fazendo o mesmo, numa escala maior. Ela tira do caminho boa parte da construção, e o
que sobra fica mais evidente do que nunca: entender o negócio, definir as regras, desenhar os
dados, conferir os resultados, decidir o que vale a pena fazer e assumir a responsabilidade pelo
que foi entregue. Não é pouco. É, na verdade, a parte do trabalho que sempre foi a mais valiosa, e
que a técnica, por ser difícil, escondia.

Há algo que atravessa todas as ondas sem mudar, e que talvez seja a principal lição deste livro.
Luhn, em 1958, já falava em levar a informação certa à pessoa certa para orientar uma ação
(\veja{sec:termo-bi}). Ackoff, em 1967, já alertava que mais informação não é melhor informação
(\veja{sec:sad}). As ferramentas mudaram tantas vezes que é difícil lembrar os nomes de todas. A
pergunta de fundo continua a mesma: que decisão este dado ajuda a tomar, e como ter certeza de que
ele está certo?

\rascunho{
  A minha aposta para os próximos anos, em primeira pessoa e assumindo que é aposta. Não repetir
  o que está na previsão sobre os times (\veja{sec:previsao-times}); pode ser sobre o estágio
  proativo da IA, sobre o papel do gerente de dados ou sobre o que eu acho que vai parecer
  ingênuo, daqui a cinco anos, em tudo o que escrevi.
}

Para quem está começando, o conselho que fecha o livro é simples, e não depende de nenhuma
ferramenta. Aprenda a técnica a fundo o bastante para saber quando a máquina errou. Aprenda um
negócio a fundo o bastante para saber o que vale a pena perguntar. Use a IA para ir mais longe, e
não para deixar de aprender. E nunca publique um número que você não saberia explicar.

\rascunho{
  Fecho pessoal opcional: uma ou duas frases minhas, de agradecimento ao leitor ou sobre o que
  me levou a escrever o livro, se não estiver no prefácio.
}
